% ============================================================
%  CHAPTERS 4–6 — PRACTICE
% ============================================================
% Answers are in chapters/answers.tex, numbered the same way.
% Forms are given for a male speaker; a female speaker uses -ii
% endings (\textit{jaatii huun}, \textit{gaii}, \textit{jaauungii}).

\section{Practice: Grammar}

Write your answers in a notebook, then check them in the Answer Key. Section numbers in brackets point to the grammar that is practised.

% ------------------------------------------------------------
\subsection{4.1 Present habitual \& continuous (1.6)}

Fill in the verb. The infinitive and the tense are given in brackets.

\begin{enumerate}
  \item \hw{मैं रोज़ चाय \blank{} हूँ।} \hfill (piinaa, habitual)
  \item \hw{मेरी बहन दिल्ली में \blank{} है।} \hfill (rahnaa, habitual)
  \item \hw{हम हिंदी \blank{} हैं।} \hfill (siikhnaa, continuous)
  \item \hw{आप क्या \blank{} हैं?} \hfill (karnaa, continuous)
  \item \hw{बच्चे पार्क में \blank{} हैं।} \hfill (khelnaa, habitual)
  \item \hw{वह} (f.) \hw{खाना \blank{} है।} \hfill (banaanaa, continuous)
\end{enumerate}

% ------------------------------------------------------------
\subsection{4.2 Simple past — \textit{ne} or not? (1.6, 1.11)}

Rewrite each sentence in the simple past. Decide whether the subject needs \textit{ne}, and make the verb agree with the right word.

\begin{enumerate}
  \item \hw{मैं ऑफिस जाता हूँ।}
  \item \hw{मैं चाय पीता हूँ।}
  \item \hw{वह} (f.) \hw{घर आती है।}
  \item \hw{हम खाना खाते हैं।}
  \item \hw{मैं किताब पढ़ता हूँ।}
  \item \hw{मैं साइकिल चलाता हूँ।}
  \item \hw{वह} (m.) \hw{काम करता है।}
\end{enumerate}

% ------------------------------------------------------------
\subsection{4.3 Future tense (1.6)}

Put the verb in the future.

\begin{enumerate}
  \item \hw{मैं कल ऑफिस \blank{}।} \hfill (jaanaa)
  \item \hw{वह} (f.) \hw{खाना \blank{}।} \hfill (banaanaa)
  \item \hw{तुम} (m.) \hw{कब \blank{}?} \hfill (aanaa)
  \item \hw{हम हिंदी \blank{}।} \hfill (siikhnaa)
  \item \hw{क्या आप पानी \blank{}?} \hfill (piinaa)
  \item \hw{वह} (m.) \hw{घर पर \blank{}।} \hfill (rahnaa)
\end{enumerate}

% ------------------------------------------------------------
\subsection{4.4 Oblique case (1.12)}

Put the word in brackets into the correct form.

\begin{enumerate}
  \item \hw{\blank{} में कौन है?} \hfill (\hw{कमरा})
  \item \hw{\blank{} को पानी दो।} \hfill (\hw{लड़का})
  \item \hw{मैं \blank{} के साथ था।} \hfill (\hw{दोस्त}, pl.)
  \item \hw{\blank{} में बहुत लोग हैं।} \hfill (\hw{बड़ा शहर})
  \item \hw{इन \blank{} में क्या है?} \hfill (\hw{किताब}, pl.)
  \item \hw{यह \blank{} के लिए है।} \hfill (\hw{लड़की}, pl.)
\end{enumerate}

% ------------------------------------------------------------
\subsection{4.5 Possessives \& \textit{apnaa} (1.13)}

Fill in the possessive. Use \textit{apnaa/apne/apnii} where the owner is the subject.

\begin{enumerate}
  \item \hw{मैं \blank{} घर जाता हूँ।} \hfill (my own)
  \item \hw{\blank{} नाम दीमा है।} \hfill (my)
  \item \hw{वह} (f.) \hw{\blank{} किताब पढ़ती है।} \hfill (her own)
  \item \hw{क्या यह \blank{} गाड़ी है?} \hfill (your, polite)
  \item \hw{\blank{} बहन मुझसे बड़ी है।} \hfill (my)
  \item \hw{हम \blank{} दोस्तों से मिलते हैं।} \hfill (our own)
\end{enumerate}

% ------------------------------------------------------------
\subsection{4.6 Imperative (1.6)}

Give the polite (\textit{aap}) and the familiar (\textit{tum}) command.

\begin{enumerate}
  \item \hw{बैठना} (sit)
  \item \hw{आना} (come)
  \item \hw{करना} (do)
  \item \hw{देना} (give)
  \item \hw{बोलना} (speak)
  \item \hw{पीना} (drink)
  \item Say ``Don't go!'' to a friend.
\end{enumerate}

% ------------------------------------------------------------
\subsection{4.7 Joining actions with \textit{-kar} (1.15)}

Join the two sentences into one. Watch out: \textit{ne} depends on the \emph{last} verb.

\begin{enumerate}
  \item \hw{मैं नाश्ता करता हूँ। मैं ऑफिस जाता हूँ।}
  \item \hw{हाथ धोइए। खाइए।}
  \item \hw{मैंने खाना खाया। मैं घर गया।}
  \item \hw{वह} (f.) \hw{घर आई। उसने चाय पी।}
\end{enumerate}

% ------------------------------------------------------------
\subsection{4.8 Subjunctive \& conditionals (1.14, 1.16)}

Translate into Hindi.

\begin{enumerate}
  \item Shall I come in?
  \item Let's go!
  \item What should I do?
  \item Maybe she will come tomorrow.
  \item If it rains, I will stay at home.
  \item If you (\textit{aap}) are tired, rest.
\end{enumerate}

% ============================================================
\clearpage
\section{Practice: Vocabulary \& Reading}

% ------------------------------------------------------------
\subsection{5.1 Masculine or feminine?}

Write m.\ or f.\ for each noun.

\medskip
\begin{tabular}{@{}llllll@{}}
  1.\ \hw{किताब} & 2.\ \hw{घर}   & 3.\ \hw{चाय}    & 4.\ \hw{पानी}   & 5.\ \hw{साइकिल} & 6.\ \hw{दिन} \\[4pt]
  7.\ \hw{रात}   & 8.\ \hw{कमरा} & 9.\ \hw{बारिश}  & 10.\ \hw{दर्द}  & 11.\ \hw{दुकान} & 12.\ \hw{फूल} \\
\end{tabular}

% ------------------------------------------------------------
\subsection{5.2 Write the Hindi word}

\begin{tabular}{@{}lll@{}}
  1.\ morning        & 5.\ left (direction)  & 9.\ yesterday / tomorrow \\
  2.\ Tuesday        & 6.\ red               & 10.\ head \\
  3.\ grandmother (father's mother) & 7.\ weather & 11.\ bicycle \\
  4.\ expensive      & 8.\ friend            & 12.\ flower \\
\end{tabular}

% ------------------------------------------------------------
\subsection{5.3 Opposites}

Write the opposite in Hindi.

\begin{tabular}{@{}llll@{}}
  1.\ \hw{बड़ा}   & 2.\ \hw{अच्छा} & 3.\ \hw{गरम}  & 4.\ \hw{महँगा} \\[4pt]
  5.\ \hw{नया}    & 6.\ \hw{दाएँ}  & 7.\ \hw{ऊपर}  & 8.\ \hw{पास} \\
\end{tabular}

% ------------------------------------------------------------
\subsection{5.4 Numbers}

Write in Hindi words: \quad 1.\ 7 \quad 2.\ 12 \quad 3.\ 15 \quad 4.\ 20 \quad 5.\ 25 \quad 6.\ 30 \quad 7.\ 50 \quad 8.\ 100

% ------------------------------------------------------------
\subsection{5.5 What time is it? (3.15)}

Answer with a full sentence (\hw{\ldots{} बजे हैं।}).

\begin{tabular}{@{}llll@{}}
  1.\ 3:00 & 2.\ 1:00 & 3.\ 4:30 & 4.\ 1:30 \\
  5.\ 2:30 & 6.\ 6:15 & 7.\ 7:45 & 8.\ 5:10 \\
\end{tabular}

% ------------------------------------------------------------
\subsection{5.6 Reading practice (1.1)}

Read aloud, then write each word in transliteration.

\medskip
\begin{tabular}{@{}llllll@{}}
  1.\ \hw{नमस्ते} & 2.\ \hw{किताब}  & 3.\ \hw{दोस्त}   & 4.\ \hw{लड़की} & 5.\ \hw{हिंदी}   & 6.\ \hw{कौन} \\[4pt]
  7.\ \hw{ठंडा}   & 8.\ \hw{कृपया}  & 9.\ \hw{स्टेशन}  & 10.\ \hw{मौसम} & 11.\ \hw{पहाड़}  & 12.\ \hw{ज़्यादा} \\
\end{tabular}

% ============================================================
\clearpage
\section{Practice: Translation}

% ------------------------------------------------------------
\subsection{6.1 Find the mistake}

Each sentence has one mistake. Write the correct sentence.

\begin{enumerate}
  \item \hw{मैं साइकिल चलाई।}
  \item \hw{मैं मेरे घर जाता हूँ।}
  \item \hw{आप क्या कर रहे हो?}
  \item \hw{वह नहीं तैर सकती।}
  \item \hw{मैंने ऑफिस गया।}
  \item \hw{लड़का को पानी दो।}
  \item \hw{मुझे चाय पसंद हूँ।}
  \item \hw{मैंने कुट्टू के फूल देखा है।}
\end{enumerate}

% ------------------------------------------------------------
\subsection{6.2 English → Hindi}

\begin{enumerate}
  \item I am learning Hindi.
  \item I have a bicycle.
  \item I need water.
  \item Please speak slowly.
  \item Yesterday I went to the office by bus.
  \item I have never eaten buckwheat.
  \item I will have to work tomorrow.
  \item My sister is older than me.
  \item If you (\textit{tum}) come, we will eat together.
  \item I have a headache.
  \item The class is at 6 pm.
  \item Today I cycled from Shakti to Agham.
\end{enumerate}

% ------------------------------------------------------------
\subsection{6.3 Hindi → English}

\begin{enumerate}
  \item \hw{मुझे आम पसंद है।}
  \item \hw{स्टेशन यहाँ से कितना दूर है?}
  \item \hw{मैंने नाश्ता नहीं किया।}
  \item \hw{वह एक गाँव में रहती थी।}
  \item \hw{मैं खाना खाकर घर गया।}
  \item \hw{क्या मैं अंदर आऊँ?}
  \item \hw{आज बहुत ठंड है।}
  \item \hw{फिर से बोलिए।}
\end{enumerate}

% ------------------------------------------------------------
\subsection{6.4 Write about yourself}

No answer key: write 4--5 sentences for each and bring them to your next lesson.

\begin{enumerate}
  \item \textbf{Yesterday} — what did you do? (simple past, \textit{ne})
  \item \textbf{This weekend} — what will you do? (future)
  \item \textbf{Your family} — who they are, where they live, what they have (\textit{ke paas}, possessives)
  \item \textbf{Your last trip} — from where to where, how long, the highest / best thing (\textit{se\ldots tak}, \textit{sabse})
  \item \textbf{Things you have never done} (\textit{kabhii nahiin} + present perfect)
\end{enumerate}
